% Vietnamese translation for OI Public Library
% Source-File: IOI中国国家候选队论文/2008/Day1/4.顾研《浅谈随机化思想在几何问题中的应用》/浅谈随机化思想在几何问题中的应用.doc
\documentclass[11pt]{article}
\usepackage{oipl-vn}

\newcommand{\Oh}{\mathcal{O}}

\title{Cảm nhận vẻ đẹp của ngẫu nhiên}
\author{Gu Yan\\Trường Trung học số 1 Trung Sơn, Quảng Đông}
\date{Trại đông Tin học toàn quốc, 2008}

\begin{document}
\maketitle

\origtitle{Cảm nhận vẻ đẹp của ngẫu nhiên: bàn về ứng dụng tư tưởng ngẫu nhiên hóa trong bài toán hình học}
\sourcepath{IOI China candidate team papers / 2008 / Day 1 / Gu Yan / randomized geometry.doc}

\renewcommand{\abstractname}{Tóm tắt}
\begin{abstract}
Trong vài năm gần đây, ngày càng nhiều bài hình học có thể giải bằng tư tưởng
ngẫu nhiên hóa. Bài viết giới thiệu hai nhóm thuật toán ngẫu nhiên thường gặp
trong thi tin học: thuật toán gia tăng ngẫu nhiên và mô phỏng tôi luyện. Qua
so sánh với các phương pháp truyền thống, bài viết nhấn mạnh ưu điểm đơn giản,
nhanh, linh hoạt và dễ song song hóa của ngẫu nhiên hóa.
\end{abstract}

\noindent\textbf{Từ khóa:} ngẫu nhiên hóa; thuật toán gia tăng ngẫu nhiên;
mô phỏng tôi luyện; điều chỉnh.

\section{Khái quát về thuật toán ngẫu nhiên}

Thuật toán ngẫu nhiên là thuật toán đưa yếu tố ngẫu nhiên vào quá trình lựa
chọn bước tiếp theo. Với cùng một đầu vào, các lần chạy khác nhau có thể có
thời gian hoặc kết quả trung gian khác nhau. Bù lại, thuật toán thường đơn
giản, nhanh, linh hoạt, và có thể cân bằng giữa thời gian, bộ nhớ và độ chính
xác.

Bản gốc chia thuật toán ngẫu nhiên thành bốn nhóm quen thuộc.

\paragraph{Thuật toán xác suất số trị.}
Loại này dùng cho bài toán số, thường cho nghiệm xấp xỉ. Khi thời gian tính
tăng, độ chính xác cũng tăng. Ví dụ kinh điển là ước lượng \(\pi\) bằng cách
lấy ngẫu nhiên điểm trong hình vuông ngoại tiếp đường tròn đơn vị.

\paragraph{Las Vegas.}
Thuật toán Las Vegas không trả lời sai: nếu tìm được đáp án thì đáp án đúng.
Tuy nhiên thời gian chạy có thể ngẫu nhiên; chạy lâu hơn thì xác suất tìm được
đáp án tăng.

\paragraph{Monte Carlo.}
Thuật toán Monte Carlo luôn trả lời trong thời gian kiểm soát được, nhưng có
thể sai với xác suất nhỏ. Nhiều phép kiểm định có lỗi một phía thuộc nhóm này,
chẳng hạn Miller--Rabin. Một nhánh khác là lấy mẫu Monte Carlo: xây một mô
hình xác suất, mô phỏng nhiều lần, rồi dùng kỳ vọng mẫu làm đáp án xấp xỉ.

\paragraph{Sherwood.}
Thuật toán Sherwood luôn cho đáp án đúng, nhưng dùng ngẫu nhiên để phá liên hệ
giữa dữ liệu xấu và hành vi xấu của thuật toán xác định. QuickSort ngẫu nhiên
là ví dụ điển hình. Thuật toán gia tăng ngẫu nhiên trong phần sau thuộc tinh
thần này.

Ngoài các nhóm trên, nhiều thuật toán điều chỉnh có yếu tố ngẫu nhiên cũng hữu
ích trong hình học: khởi tạo một phương án, rồi liên tục điều chỉnh tới khi đạt
điều kiện mong muốn.

\section{Thuật toán gia tăng ngẫu nhiên}

\subsection{Từ gia tăng tới gia tăng ngẫu nhiên}

Thuật toán gia tăng xử lý bài toán bằng cách giải bài con nhỏ hơn rồi thêm đối
tượng mới vào. Nếu \(S_i\) là tập gồm \(i\) đối tượng đầu tiên, ta duy trì đáp
án của \(S_i\), sau đó thêm đối tượng thứ \(i+1\) để thu được đáp án của
\(S_{i+1}\).

Cách này thường dễ cài đặt nhưng có thể chậm trong trường hợp xấu, vì thứ tự
đầu vào có thể ép thuật toán sửa đáp án quá nhiều lần. Ý tưởng Sherwood là xáo
trộn ngẫu nhiên thứ tự đầu vào trước khi chạy. Khi đó thời gian kỳ vọng được
tính theo thứ tự ngẫu nhiên, không phụ thuộc vào việc dữ liệu gốc được sắp xếp
xấu ra sao.

Một phép xáo trộn tuyến tính tiêu chuẩn là Fisher--Yates:
\begin{verbatim}
for i = 2..n:
    swap(a[i], a[random(1..i)])
\end{verbatim}

\subsection{Ví dụ: đường tròn bao nhỏ nhất}

Cho \(n\) điểm trên mặt phẳng, cần tìm đường tròn nhỏ nhất chứa tất cả điểm.
Một thuật toán xác định dựa trên việc xét các cặp hoặc bộ ba điểm có thể dẫn
tới độ phức tạp cao. Thuật toán gia tăng ngẫu nhiên dựa trên tính chất sau:
nếu \(D(P)\) là đường tròn bao nhỏ nhất của tập \(P\), và thêm một điểm \(p\):
\begin{itemize}
  \item nếu \(p\in D(P)\), đáp án không đổi;
  \item nếu \(p\notin D(P)\), thì \(p\) nằm trên biên của đường tròn bao mới.
\end{itemize}

Vì vậy sau khi xáo trộn các điểm, ta lần lượt thêm từng điểm. Nếu điểm mới còn
nằm trong đường tròn hiện tại thì bỏ qua. Nếu không, ta dựng lại đường tròn
với điểm mới buộc nằm trên biên, bằng cách xét các điểm trước đó. Có thể tiếp
tục dùng cùng ý tưởng: khi thêm điểm thứ hai làm vi phạm, hai điểm ấy cùng nằm
trên biên; khi thêm điểm thứ ba làm vi phạm, ba điểm xác định đường tròn.

Phân tích kỳ vọng dựa trên tuyến tính của kỳ vọng. Ở bước \(i\), xác suất điểm
thứ \(i\) nằm trên biên đường tròn bao của \(i\) điểm chỉ cỡ \(\Oh(1/i)\), vì
đường tròn được xác định bởi nhiều nhất ba điểm biên. Do đó số lần phải dựng
lại ở các bước lớn là nhỏ theo kỳ vọng. Thuật toán đạt thời gian tuyến tính kỳ
vọng nếu cài đặt đầy đủ; ngay cả phiên bản đơn giản hơn cũng đã tốt hơn nhiều
so với cách xét mọi bộ điểm.

Điểm quan trọng là chữ ``kỳ vọng'' ở đây xét theo thứ tự xáo trộn ngẫu nhiên,
không phải theo phân bố đầu vào.

\subsection{Ví dụ: Expensive Drink}

Bài \textit{Expensive Drink} có thể chuyển thành quy hoạch tuyến tính ba chiều.
Gọi \(x,y,z\) lần lượt là giá của nước, cà phê và rượu. Mỗi đồ uống đã biết
tạo ra hai ràng buộc tuyến tính từ khoảng giá phần đường, kèm các ràng buộc
\[
  0\le x\le y\le z.
\]
Cần cực đại hóa một hàm tuyến tính của \(x,y,z\).

Cách thô là lấy giao của mọi bộ ba mặt phẳng, kiểm tra tất cả ràng buộc, rồi
cập nhật đáp án. Nếu có \(m\) mặt phẳng, chi phí có thể tới \(\Oh(m^4)\).

Cải tiến đầu tiên là xét giao của hai mặt phẳng, tức một đường thẳng. Các ràng
buộc còn lại cắt đường thẳng này thành một đoạn hợp lệ. Vì hàm mục tiêu tuyến
tính trên một đoạn đạt cực trị ở một đầu mút, chỉ cần kiểm tra hai đầu đoạn.
Cách này giảm chi phí nhưng vẫn chưa đủ đẹp.

Thuật toán gia tăng xử lý từng nửa không gian. Giả sử sau \(i\) ràng buộc ta
đã có nghiệm tối ưu \(v_i\) của phần giao hiện tại. Khi thêm ràng buộc
\(h_{i+1}\):
\begin{itemize}
  \item nếu \(v_i\) vẫn thỏa \(h_{i+1}\), đáp án không đổi;
  \item nếu không, nghiệm tối ưu mới nếu tồn tại phải nằm trên mặt phẳng biên
  của \(h_{i+1}\).
\end{itemize}
Khi đó bài toán giảm xuống bài trên mặt phẳng, có thể giải bằng phiên bản cắt
đoạn. Nếu trước khi xử lý ta xáo trộn thứ tự ràng buộc, xác suất một ràng buộc
mới thật sự xác định nghiệm tối ưu ở bước \(i\) chỉ cỡ \(\Oh(1/i)\). Tổng thời
gian kỳ vọng vì vậy giảm mạnh so với các cách liệt kê trực tiếp.

Bản gốc nhấn mạnh rằng, mặc dù tồn tại các thuật toán tuyến tính lý thuyết cho
quy hoạch tuyến tính số chiều cố định, hằng số và độ phức tạp cài đặt của
chúng không phù hợp với nhiều cuộc thi. Thuật toán gia tăng ngẫu nhiên đạt cân
bằng tốt hơn giữa ý tưởng, mã nguồn và hiệu năng.

\section{Mô phỏng tôi luyện}

\subsection{Mô hình chung}

Mô phỏng tôi luyện là một meta-heuristic lấy cảm hứng từ quá trình làm nguội
vật rắn. Thuật toán có trạng thái hiện tại \(S\), nhiệt độ \(T\), và hàm đánh
giá \(C(S)\). Ở mỗi bước sinh trạng thái mới \(S'\), tính
\[
  \Delta=C(S')-C(S).
\]
Nếu \(S'\) tốt hơn, nhận nó. Nếu xấu hơn, vẫn có thể nhận với xác suất xấp xỉ
\[
  \exp\!\left(-\frac{\Delta}{T}\right).
\]
Khi \(T\) giảm dần, thuật toán chuyển từ khám phá rộng sang tinh chỉnh cục bộ.

Trong hình học, mô phỏng tôi luyện đặc biệt hợp với các bài tối ưu vị trí mà
hàm mục tiêu biến đổi liên tục, ví dụ cực đại hóa khoảng cách tới điểm gần
nhất, cực tiểu hóa tổng khoảng cách, hoặc tìm vị trí tốt theo một tiêu chuẩn
khó rời rạc hóa.

\subsection{Run Away}

Bài \textit{Run Away} cho một hình chữ nhật và nhiều lỗ nguy hiểm. Cần tìm
điểm trong hình chữ nhật sao cho khoảng cách tới lỗ gần nhất là lớn nhất.

Lời giải hình học chính xác liên quan tới Voronoi diagram: ứng viên tối ưu
nằm ở đỉnh Voronoi hoặc giao của cạnh Voronoi với biên hình chữ nhật. Nhưng
cài đặt Fortune, Delaunay hoặc chia để trị nửa mặt phẳng có thể rất dài.

Mô phỏng tôi luyện dùng một cách nhìn khác:
\begin{itemize}
  \item trạng thái là một điểm trong hình chữ nhật;
  \item hàm đánh giá là khoảng cách tới lỗ gần nhất;
  \item ở mỗi nhiệt độ, sinh điểm mới bằng cách đi một đoạn ngẫu nhiên độ dài
  xấp xỉ \(\Delta\);
  \item nếu điểm mới tốt hơn, giữ nó; nếu không, có thể bỏ qua hoặc nhận theo
  xác suất;
  \item giảm \(\Delta\) cho tới khi đủ nhỏ.
\end{itemize}

Để tránh phụ thuộc vào một điểm khởi đầu, có thể duy trì nhiều ứng viên cùng
lúc rồi lấy điểm tốt nhất. Thuật toán ngắn, dễ điều chỉnh thời gian chạy, và
thường đủ tốt khi yêu cầu độ chính xác không quá khắc nghiệt.

\subsection{Empire Strikes Back}

Bài này yêu cầu bán kính bom nhỏ nhất sao cho các bom đặt tại các nhà máy phủ
được toàn bộ đĩa quốc gia. Đây cũng là bài tối ưu khoảng cách: cần cực tiểu
hóa giá trị lớn nhất của khoảng cách từ một điểm chưa phủ tới nhà máy gần
nhất, theo cách diễn đạt tương đương của vùng cần bao phủ.

Bản gốc dùng cùng tư tưởng mô phỏng tôi luyện: chọn các điểm ứng viên, đánh
giá mức ``xấu nhất'', rồi dần dịch chuyển để cải thiện hàm mục tiêu. So với
việc xây dựng cấu trúc hình học chính xác, cách này có ưu thế ở mã ngắn và dễ
thay đổi tham số.

\subsection{Laser Tank}

\textit{Laser Tank} là bài nộp đáp án. Cần chọn vị trí đặt bộ phát laser sao
cho bắn tới các mục tiêu với số lần phản xạ bị giới hạn, tối ưu theo hàm điểm
của đề. Với \(k=0\), ta chỉ kiểm tra đường thẳng không bị gương cản. Với
\(k=1\), có thể xét thêm đường phản xạ qua từng gương. Với \(k=2\) trở lên,
số trường hợp tăng nhanh và cài đặt trở nên phức tạp.

Mô phỏng tôi luyện hữu ích vì ta có thể kiểm soát thời gian: trước hết tìm một
vùng tốt bằng ít vòng lặp, sau đó tinh chỉnh quanh vùng ấy. Khi cần thêm trường
hợp phản xạ, ta thay hàm đánh giá mà vẫn giữ khung tìm kiếm. Trong bài nộp đáp
án, một lời giải ngắn cho điểm khá thường có giá trị thực tế hơn một lời giải
chính xác nhưng quá dài để viết trong thời gian thi.

\section{Kết luận}

Ngẫu nhiên hóa không thay thế nền tảng hình học cổ điển. Ngược lại, muốn dùng
ngẫu nhiên hóa hiệu quả vẫn phải hiểu rõ cấu trúc bài toán: điểm nào xác định
đường tròn, mặt phẳng nào quyết định nghiệm tối ưu, hay hàm mục tiêu hình học
thay đổi ra sao. Nhưng khi đã hiểu cấu trúc ấy, ngẫu nhiên hóa thường đem lại
một đường đi nhẹ hơn: mã ngắn hơn, phân tích trung bình tốt hơn, và dễ cân đối
giữa thời gian với chất lượng nghiệm.

Thuật toán gia tăng ngẫu nhiên cho thấy cách biến một thuật toán dễ viết nhưng
nhạy với thứ tự đầu vào thành thuật toán có kỳ vọng tốt. Mô phỏng tôi luyện
cho thấy cách tìm nghiệm gần tối ưu trong các bài hình học liên tục mà cấu
trúc chính xác quá nặng. Đó chính là vẻ đẹp của ngẫu nhiên trong thi tin học.

\section*{Lời cảm ơn}

Tác giả cảm ơn huấn luyện viên Liu Rujia, bạn Yu Huacheng, nhóm tin học của
Trường Trung học số 1 Trung Sơn, và các tài nguyên trên Internet đã hỗ trợ quá
trình nghiên cứu.

\begin{thebibliography}{9}
\bibitem{manber}
Udi Manber, \textit{Introduction to Algorithms: A Creative Approach}.

\bibitem{alsuwaiyel}
M. H. Alsuwaiyel, \textit{Algorithms: Design Techniques and Analysis}.

\bibitem{algorithm-art}
Liu Rujia và Huang Liang, \textit{Nghệ thuật thuật toán và thi tin học}.

\bibitem{cg}
M. de Berg và cộng sự, \textit{Computational Geometry: Algorithms and
Applications}.

\bibitem{tang}
Tang Wenbin, \textit{Bàn về tư tưởng điều chỉnh trong thi tin học}, luận văn
đội tuyển quốc gia 2006.

\bibitem{wang}
Wang Dong, \textit{Bàn về xây dựng và ứng dụng Voronoi diagram phẳng}, luận
văn đội tuyển quốc gia 2007.
\end{thebibliography}

\section*{Ghi chú về phụ lục}

Bản gốc kèm nhiều hình minh họa, chứng minh bổ sung, phát biểu đề và các đoạn
giả mã/chương trình. Các hình được diễn giải bằng lời; các chương trình dài và
phát biểu đầy đủ của đề không được chép lại. Những giả mã ngắn có vai trò
trực tiếp trong lập luận đã được mô tả trong phần nội dung chính.

\end{document}
